\documentclass[10pt,a4paper]{amsart}
\usepackage{amsmath,amssymb,amsthm,graphicx}
\usepackage{amsmath}
\usepackage{amsfonts}
\usepackage{amssymb}
\usepackage{graphicx}
\usepackage{float}
\usepackage{caption}
\usepackage{pgf,tikz}
\usetikzlibrary{arrows}
\usepackage{verbatim}
\usepackage{hyperref}

\definecolor{ffttww}{rgb}{1,0.2,0.4}
\definecolor{uuuuuu}{rgb}{0.27,0.27,0.27}
\definecolor{cqcqcq}{rgb}{0.75,0.75,0.75}
\definecolor{wwqqcc}{rgb}{0.4,0,0.8}

\newcommand{\R}{\mathbb{R}}
\newcommand{\N}{\mathbb{N}}
\newcommand{\HH}{\mathcal{H}}
\newcommand{\RR}{\mathcal{R}}
\newcommand{\A}{\mathcal{A}}
\newcommand{\B}{\mathcal{B}}
\newcommand{\est}{\emptyset}
\newtheorem{theorem}{Theorem}
\newtheorem{proposition}{Proposition}
\newtheorem{corollary}{Corollary}
\newtheorem{lemma}{Lemma}
\newtheorem{definition}{Definition}
\newtheorem{jel}{Notation}
\newtheorem{conj}{Conjecture}
\title{Strict lower bound on the polychromatic parameter of dual axis-parallel strip hypergraphs}
\author{Balázs Bursics, Bence Csonka, Luca Szepessy}
\date{arxiv:2307.12154v3; CC BY 4.0}
\begin{document}
\maketitle
\noindent\textit{Editorial note.} The counted statement is printed in the introduction of the article, while its proof stands in the Proofs section; the proof environment carries the author's own attribution, Proof of Theorem strip\_dual, which is the sole pairing used here. No other proof of this theorem exists in the article.\par\medskip
\noindent\textit{Editorial note.} The article's two transported PNG figures (abra1\_en.png, abra3.png) belong to the separate strip construction and are not part of this item, so no asset is declared and no figure is deleted from the spans that are declared.\par\medskip
\noindent\textit{Editorial note.} The retained preamble is the native macro and theorem-environment block of the article; font, spacing, caption and penguin-drawing packages of the original file are dropped, and the animation-free statement, proof and context spans need none of them.\par\medskip
\noindent\textit{Editorial note.} Only the counted theorem is a counting unit. The article's Theorem strip, Theorem kereszt, Theorem s\_strip, Theorem vegtelen\_sorozat, Theorem veges\_sorozat and the corollary printed after s\_strip are uncounted neighbours of this card.\par\medskip
\noindent\textit{Editorial scope.} One counted lower bound for the dual axis-parallel strip family: 2(ceil(3k/4)-1)+1 <= m\_\{A*\}(k), with the article's complete four-family counting proof. The primal strip and bottomless non-existence theorems, the A-plus and A\_s bounds, both arithmetic-progression theorems and the printed corollary after s\_strip are not counted.\par\medskip
\noindent\textit{Allowed standard prerequisites (editorial).}\begin{itemize}
\item Polychromatic colourings and the parameter m\_H(k). A polychromatic k-colouring of a hypergraph puts a vertex of every one of the k colour classes into every hyperedge; m\_H(k) is the least m such that every member of the family has a polychromatic k-colouring after all edges of size below m are deleted, and it is infinity when no such m exists.. source lines 120 and 123-129, retained as context. Supplies the counted quantity m\_\{A*\}(k) and the meaning of an edge that is not polychromatic in the dual family.
\item Finite plane point sets, strips and slight perturbation. A finite vertex set in the plane and finitely many axis-parallel strips realise the dual family A-star, whose vertices are the strips and whose edges are the sets of strips containing one point. Finitely many parallel strips can be perturbed a little so that no two of them coincide, while every nonempty region of their common intersection still captures the same families.. parenthetical sentence in the counted proof, source line 479. The perturbation sentence at source line 479 keeps the four families of copies distinct as vertices of A-star and keeps the private part of each pair intersection captured by exactly the two corresponding families.
\item Elementary finite counting and the pigeonhole principle. A set of c vertices carries at most c colours, so at least k minus c of the k colours are absent from it; four subsets of a k-element set whose sizes sum to more than k cannot be pairwise disjoint.. source line 481. Gives the more-than-k/4 missing colours per family and the shared missing colour of two families in the counted proof.
\end{itemize}\medskip
A {\it polychromatic k-coloring} of a hypergraph $H$ is a $k$-coloring of its vertex set such that every hyperedge contains a vertex of all $k$ color classes. Observe that a polychromatic 2-coloring is the same as the usual proper 2-coloring of hypergraphs, where we require that no edge is monochromatic. By merging $j$ color classes of a polychromatic $k$-coloring we get a polychromatic $(k-j+1)$-coloring, so the condition of monochromatic $k$-colorability becomes stricter as $k$ increases. A trivial necessary condition for the existence of a polychromatic $k$-coloring is that all edges of $H$ must be of size at least $k$.

\begin{definition} Denote by $H_{\ge m}$ or simply $H_m$ the hypergraph obtained from $H$ by deleting all hyperedges of size smaller than $m$, and denote by $H_{=m}$ the hypergraph consisting of the hyperedges of $H$ with size exactly $m$. Similarly, for a hypergraph family $\mathcal{H}$ let $\HH_m=\HH_{\ge m}=\{H_{\ge m}:H\in\HH\}$ and $\HH_{=m}=\{H_{=m}:H\in\HH\}$.
\end{definition}

\begin{definition} Let $\mathcal{H}$ be a hypergraph family. Denote by $m_\mathcal{H}(k)$ the smallest positive integer $m$ such that for every $H\in\mathcal{H}$ there exists a polychromatic $k$-coloring of the hypergraph $H_{\ge m}$. If there is no such $m$, set $m_{\mathcal{H}}(k)=\infty$.
\end{definition}

\begin{definition} Denote by $\A$ the hypergraph family which consists of those hypergraphs $H=(V,E)$ for which $V\subset \R^2$ is finite, and the edges are induced by axis-parallel strips: every $e\in E$ edge can be presented in the form $V\cap\{(x,y)\in\R^2: x_0<x<x_1\}$ or $V\cap\{(x,y)\in\R^2: y_0<y<y_1\}$ for some $x_0, x_1, y_0, y_1 \in\R$. Denote by $\A^*$ its dual: the family of such hypergraphs that the vertices are axis-parallel strips, and the edges are sets consisting of the strips containing a given point $(x,y)$. 
\end{definition}

\begin{theorem}\label{strip_dual}
For the hypergraph family $\A^*$ we have 
$2\left(\left\lceil \frac{3}{4}k\right\rceil -1\right) + 1 \leq m_{\A^*}(k)$.
\end{theorem}

\begin{proof}[Proof of Theorem~\ref{strip_dual}]

Take $\lceil \frac{3}{4}k\rceil-1$ copies each of the following axis-parallel strips: $\{(x,y):0<x<2\}$, $\{(x,y):1<x<3\}$, $\{(x,y):0<y<2\}$, and $\{(x,y):1<y<3\}$. Notice that any two of these has a part of their intersection which is disjoint from the other two original strips. (If we do not want vertices derived from the same strip, we can perturb them a little.)

Take that $H$ hypergraph from $\A^*$ which has these $4(\lceil \frac{3}{4}k\rceil-1)$ strips as vertices and all possible hyperedges determined by points. In any coloring of the strips, for each of the 4 original strips there exist strictly more than $\frac{k}{4}$ colors which do not appear in the copies of that strip. By the pigeonhole principle there are two original strips and a color excluded from copies of these two strips, hence $H$ has a hyperedge of size $2(\lceil \frac{3}{4}k\rceil-1)$ which is not polychromatic, implying $2\left(\left\lceil \frac{3}{4}k\right\rceil -1\right) + 1 \leq m_{\A^*}(k)$.
\end{proof}

\end{document}
